% !TEX program = lualatex
\documentclass[a4paper,11pt]{jlreq}

\usepackage{luatexja-fontspec}
\usepackage{graphicx}
\usepackage{xcolor}
\usepackage[most]{tcolorbox}
\usepackage{tikz}
\usetikzlibrary{positioning,arrows.meta,shapes.geometric,shapes.symbols,fit,backgrounds,calc,decorations.pathreplacing,patterns,shadows.blur}
\usepackage{enumitem}
\usepackage{array}
\usepackage{tabularx}
\usepackage{booktabs}
\usepackage{titlesec}
\usepackage{fontawesome5}
\usepackage[hidelinks]{hyperref}
\usepackage{needspace}
\usepackage{geometry}
\geometry{a4paper,top=22mm,bottom=22mm,left=20mm,right=20mm}

% --- カラーパレット ---
\definecolor{ipBlue}{HTML}{4F86F7}
\definecolor{ipBlueD}{HTML}{2A5BD7}
\definecolor{ipNavy}{HTML}{1E3A5F}
\definecolor{ipGreen}{HTML}{2DB67C}
\definecolor{ipOrange}{HTML}{F2994A}
\definecolor{ipPink}{HTML}{EC6F9C}
\definecolor{ipPurple}{HTML}{9B7BE0}
\definecolor{ipRed}{HTML}{E26565}
\definecolor{ipGray}{HTML}{6B7280}
\definecolor{ipGrayL}{HTML}{E5E9F2}
\definecolor{ipBg}{HTML}{F6F8FD}
\definecolor{ipSoft}{HTML}{FFF8EE}
\definecolor{ipMint}{HTML}{EAF8F2}
\definecolor{ipPanel}{HTML}{FFFFFF}

\titleformat{\section}{\sffamily\bfseries\LARGE\color{ipNavy}}{}{0pt}{}
\titleformat{\subsection}{\sffamily\bfseries\large\color{ipNavy}}{}{0pt}{}

\setlength{\parindent}{0pt}
\setlength{\parskip}{4pt}

% --- 章ヘッダ ---
\newcommand{\chap}[2]{%
  \par\vspace{0.4em}\noindent
  \begin{tikzpicture}
    \node[fill=ipBlue,text=white,rounded corners=4pt,minimum height=2.0em,inner xsep=10pt,inner ysep=4pt,font=\sffamily\bfseries\large] (n1) {0#1};
    \node[anchor=west,font=\sffamily\bfseries\LARGE,text=ipNavy] at ([xshift=10pt]n1.east) {#2};
  \end{tikzpicture}\par
  \vspace{-0.05em}\textcolor{ipBlue!50}{\rule{\linewidth}{1.4pt}}\vspace{-0.1em}\textcolor{ipBlue!20}{\rule{\linewidth}{0.4pt}}\par\vspace{0.4em}
}

\newtcolorbox{screenbox}[2][]{%
  enhanced, breakable,
  colback=ipPanel, colframe=ipBlue,
  boxrule=0.8pt, arc=4pt,
  left=12pt, right=12pt, top=14pt, bottom=10pt,
  fonttitle=\sffamily\bfseries\large,
  coltitle=white,
  colbacktitle=ipBlue,
  attach boxed title to top left={xshift=14pt,yshift=-12pt},
  boxed title style={enhanced,arc=4pt,colback=ipBlue,boxrule=0pt,top=3pt,bottom=3pt,left=10pt,right=10pt,interior style={fill=ipBlue}},
  drop fuzzy shadow=ipNavy!20,
  title={#2}, #1
}

\newtcolorbox{actionbox}[1][]{%
  enhanced, breakable,
  colback=ipBg, colframe=ipBlue!12,
  boxrule=0pt,
  left=12pt, right=12pt, top=7pt, bottom=7pt,
  arc=3pt,
  borderline west={2.4pt}{0pt}{ipBlue!60},
  #1
}

\newtcolorbox{tipbox}[1][]{%
  enhanced, breakable,
  colback=ipSoft, colframe=ipOrange!50,
  boxrule=0.6pt,
  left=12pt, right=12pt, top=8pt, bottom=8pt,
  arc=3pt,
  borderline west={2.4pt}{0pt}{ipOrange},
  #1
}

\newcommand{\btn}[1]{\tikz[baseline=(b.base)]{\node[fill=ipBlue,text=white,rounded corners=4pt,inner xsep=6pt,inner ysep=2.5pt,font=\sffamily\small\bfseries](b){#1};}}
\newcommand{\btnG}[1]{\tikz[baseline=(b.base)]{\node[fill=ipGreen,text=white,rounded corners=4pt,inner xsep=6pt,inner ysep=2.5pt,font=\sffamily\small\bfseries](b){#1};}}
\newcommand{\btnO}[1]{\tikz[baseline=(b.base)]{\node[fill=ipOrange,text=white,rounded corners=4pt,inner xsep=6pt,inner ysep=2.5pt,font=\sffamily\small\bfseries](b){#1};}}
\newcommand{\btnR}[1]{\tikz[baseline=(b.base)]{\node[fill=ipRed,text=white,rounded corners=4pt,inner xsep=6pt,inner ysep=2.5pt,font=\sffamily\small\bfseries](b){#1};}}
\newcommand{\tab}[1]{\tikz[baseline=(b.base)]{\node[fill=ipNavy!8,text=ipNavy,rounded corners=2.5pt,inner xsep=5pt,inner ysep=2pt,font=\sffamily\small,draw=ipNavy!25](b){#1};}}
\newcommand{\arrowto}{\,\textcolor{ipBlue}{$\Rightarrow$}\,}

\newcommand{\hl}[1]{\textbf{\textcolor{ipNavy}{#1}}}

\newcommand{\progbar}[2]{%
  \begin{tikzpicture}[baseline=-1.6pt]
    \fill[ipGrayL,rounded corners=2pt] (0,0) rectangle (2.6,0.26);
    \fill[#2,rounded corners=2pt] (0,0) rectangle (#1*0.026,0.26);
  \end{tikzpicture}%
}

% ---------------------------
\begin{document}

% =========================================
% 表紙
% =========================================
\thispagestyle{empty}
\newgeometry{margin=0pt}
\begin{tikzpicture}[remember picture,overlay]
  \fill[ipBlue] (current page.north west) rectangle ($(current page.north east)+(0,-12)$);
  \fill[ipBlueD,opacity=0.55] (current page.north west) rectangle ($(current page.north east)+(0,-12)$);
  \fill[white,opacity=0.07] ($(current page.north east)+(-3.5,-0.8)$) circle (3.4);
  \fill[white,opacity=0.05] ($(current page.north west)+(1.2,-10.5)$) circle (2.4);
  \fill[white,opacity=0.05] ($(current page.north east)+(-1.5,-10)$) circle (1.6);
  \fill[white,opacity=0.04] ($(current page.north west)+(8,-2.5)$) circle (1.2);

  \node[anchor=north west,text=white,font=\sffamily\bfseries\fontsize{15}{18}\selectfont] at ($(current page.north west)+(2.0,-1.7)$) {iPlus Sys};
  \node[anchor=north west,text=white!85,font=\sffamily\fontsize{10}{12}\selectfont] at ($(current page.north west)+(2.0,-2.3)$) {塾の運用支援アプリ};

  \node[anchor=north west,text=white,font=\sffamily\bfseries\fontsize{32}{40}\selectfont] at ($(current page.north west)+(2.0,-5.4)$) {点数を入れるだけで、};
  \node[anchor=north west,text=white!92,font=\sffamily\bfseries\fontsize{22}{28}\selectfont] at ($(current page.north west)+(2.0,-7.6)$) {進捗が進み、};
  \node[anchor=north west,text=white!92,font=\sffamily\bfseries\fontsize{22}{28}\selectfont] at ($(current page.north west)+(2.0,-8.7)$) {合格力が見えてくる。};

  \node[anchor=north west,fill=white,rounded corners=4pt,text width=10cm,inner sep=14pt,drop fuzzy shadow=ipNavy!25,
        font=\sffamily] at ($(current page.north west)+(2.0,-13.5)$)
    {\textbf{\textcolor{ipNavy}{\large ご報告資料}}\\[4pt]
     \textcolor{ipGray}{現状の機能・使い方・これから ── 塾長へ}\\[6pt]
     \textcolor{ipGray}{2026年4月28日　／　森 祐太}};

  \node[anchor=north west,fill=white,draw=ipBlue!25,line width=0.6pt,rounded corners=5pt,
        minimum width=5.0cm,minimum height=2.6cm,align=center,inner sep=8pt,drop fuzzy shadow=ipNavy!12,
        font=\sffamily] (c1) at ($(current page.north west)+(2.0,-19.5)$)
    {\tikz\node[fill=ipBlue!12,text=ipBlue,circle,minimum size=1.9em,font=\bfseries](){\faExclamationTriangle};\\[3pt]
     \bfseries\small\textcolor{ipNavy}{要注意の自動検出}\\[1pt]
     \scriptsize\textcolor{ipGray}{完了間近・正答率低下を発見}};
  \node[anchor=north west,fill=white,draw=ipGreen!30,line width=0.6pt,rounded corners=5pt,
        minimum width=5.0cm,minimum height=2.6cm,align=center,inner sep=8pt,drop fuzzy shadow=ipNavy!12,
        font=\sffamily] (c2) at ($(c1.north east)+(0.3,0)$)
    {\tikz\node[fill=ipGreen!12,text=ipGreen,circle,minimum size=1.9em,font=\bfseries](){\faClipboardCheck};\\[3pt]
     \bfseries\small\textcolor{ipNavy}{点数入力で進捗自動進行}\\[1pt]
     \scriptsize\textcolor{ipGray}{合格チェックで次の章へ}};
  \node[anchor=north west,fill=white,draw=ipPurple!30,line width=0.6pt,rounded corners=5pt,
        minimum width=5.0cm,minimum height=2.6cm,align=center,inner sep=8pt,drop fuzzy shadow=ipNavy!12,
        font=\sffamily] (c3) at ($(c2.north east)+(0.3,0)$)
    {\tikz\node[fill=ipPurple!12,text=ipPurple,circle,minimum size=1.9em,font=\bfseries](){\faCalculator};\\[3pt]
     \bfseries\small\textcolor{ipNavy}{圧縮スコアで合格力把握}\\[1pt]
     \scriptsize\textcolor{ipGray}{志望大学の配点で換算}};

  \node[anchor=south,text=ipGray,font=\sffamily\scriptsize] at ($(current page.south)+(0,1.3)$)
    {\faBookOpen\ \  全 10 ページ\quad/\quad 読み物として10分くらい};
\end{tikzpicture}
\clearpage
\restoregeometry

% =========================================
% TL;DR
% =========================================
\begin{tcolorbox}[colback=ipBg,colframe=ipBlue,boxrule=0pt,arc=5pt,left=16pt,right=16pt,top=12pt,bottom=12pt]
\sffamily
{\fontsize{14}{18}\selectfont\textbf{\textcolor{ipBlue}{\faBolt}\ \ ひとことで言うと}}\\[6pt]
{\fontsize{12}{18}\selectfont\textcolor{ipNavy}{
\textbf{教材を作る → 生徒を選ぶ → 割り当てる → 結果を入れる → 反映 → 個別印刷}。\\
この6ステップが、すべて1つのアプリの中で完結します。
採点結果を入れた瞬間に進捗が更新され、要注意の生徒には自動で印が立ちます。}}
\end{tcolorbox}

\vspace{0.2em}

\begin{tipbox}
\sffamily
\textbf{\textcolor{ipOrange}{\faLightbulb}\ \ この資料の読み方}\\[2pt]
パソコンが苦手な方でも分かるように、専門用語はなるべく使わずに書きました。
画面ごとに「\textbf{ここを押すと、こうなる}」をボックスで並べてあります。
\textbf{気になるところだけ拾い読み}でも大丈夫です。
\end{tipbox}

\vspace{0.4em}

% =========================================
% 章1 全体イメージ
% =========================================
\chap{1}{まずは全体イメージ}

ブラウザで使う、\hl{塾の管理画面}です。アプリのインストールは要りません。
データはインターネット上に保管され、教室のパソコン・iPad・スマホ、どれからでも同じ画面が開けます。

\vspace{0.4em}
\begin{center}
\begin{tikzpicture}[
  font=\sffamily\small,
  box/.style={rounded corners=4pt,draw=ipNavy!30,fill=white,minimum width=3.6cm,minimum height=1.4cm,align=center,inner sep=4pt,line width=0.7pt,drop fuzzy shadow=ipNavy!12},
  prn/.style={rounded corners=3pt,draw=ipNavy!30,fill=white,minimum width=2.6cm,minimum height=1.0cm,align=center,inner sep=3pt,font=\sffamily\small,line width=0.7pt,drop fuzzy shadow=ipNavy!12},
  arr/.style={-{Stealth[length=2.6mm]},line width=1pt,ipBlue!70},
  lab/.style={font=\sffamily\scriptsize\bfseries,fill=ipBg,inner sep=1.4pt,rounded corners=1.5pt,text=ipNavy},
  ic/.style={font=\sffamily\bfseries\large},
]
\node[box,top color=ipBlue!8,bottom color=ipBlue!18] (fe) {\tikz{\node[ic,text=ipBlue]{\faDesktop};}\\[1pt]\textbf{管理画面}\\ \scriptsize（塾長・先生が見る）};
\node[box,top color=ipGreen!8,bottom color=ipGreen!18,right=2.0cm of fe] (be) {\tikz{\node[ic,text=ipGreen]{\faDatabase};}\\[1pt]\textbf{データ保管庫}\\ \scriptsize（インターネット上）};
\node[box,top color=ipOrange!10,bottom color=ipOrange!22,right=2.0cm of be] (ag) {\tikz{\node[ic,text=ipOrange]{\faCogs};}\\[1pt]\textbf{印刷役プログラム}\\ \scriptsize（教室のパソコン）};
\node[prn,top color=ipNavy!4,bottom color=ipNavy!12,below=1.0cm of ag] (printer) {\tikz{\node[ic,text=ipNavy!70]{\faPrint};}\\[1pt]プリンタ};

\draw[arr] (fe) -- node[lab]{操作} (be);
\draw[arr] (be) -- node[lab]{印刷依頼} (ag);
\draw[arr] (ag) -- node[lab]{紙が出る} (printer);
\end{tikzpicture}
\end{center}

\begin{itemize}[leftmargin=1.4em,itemsep=3pt]
  \item \hl{管理画面}：先生・塾長が見る部分。\textbf{すべてここで完結}します。
  \item \hl{データ保管庫}：インターネット上にあるので、教室が停電しても消えません。
  \item \hl{印刷役プログラム}：教室のプリンタにつながったパソコンで動く小さなソフト。管理画面からの印刷依頼を受け取って\textbf{勝手に紙を出して}くれます。
\end{itemize}

% =========================================
% 章2 画面マップ
% =========================================
\chap{2}{どんな画面があるか}

画面は全部で9種類。普段よく使うのは上の3〜4個です。

\vspace{0.4em}
\begin{center}\sffamily\small
\renewcommand{\arraystretch}{1.45}
\begin{tabularx}{\linewidth}{@{}c >{\bfseries\color{ipNavy}}l X@{}}
\toprule
& 画面 & どんな画面か \\
\midrule
\faTachometerAlt\,\textcolor{ipBlue}{} & ダッシュボード & 塾全体の様子＋\textbf{要注意の生徒に自動で印} \\
\faUserGraduate\,\textcolor{ipBlue}{} & 生徒 & 生徒の追加／教材の割り当て／\textbf{定着度入力}／分析 \\
\faBook\,\textcolor{ipBlue}{} & 教材管理 & 教材を1冊ずつ登録（章ごと、PDFを4種類紐付け可） \\
\faSpellCheck\,\textcolor{ipBlue}{} & 単語テスト & 単語帳をエクセル感覚で\textbf{まとめて取り込み} \\
\faChartBar\,\textcolor{ipBlue}{} & 試験管理 & 共テ・過去問の管理／成績入力／\textbf{8種のグラフで分析} \\
\faPrint\,\textcolor{ipBlue}{} & 印刷 & 生徒ごとに「印刷」ボタン／問題と解答を1冊に結合 \\
\faFilePdf\,\textcolor{ipBlue}{} & PDF置き場 & 教材PDFのアップロードと一覧 \\
\faChalkboardTeacher\,\textcolor{ipBlue}{} & 講師管理 & 講師の名前を登録（定着度入力時に紐付け） \\
\faUserShield\,\textcolor{ipBlue}{} & アカウント管理 & 誰がこの画面に入れるか・権限の設定 \\
\bottomrule
\end{tabularx}
\end{center}

\vspace{0.3em}
\textcolor{ipGray}{\small※ 普段使うのは「ダッシュボード」「生徒」「印刷」が中心。教材登録などは初期設定のときだけ触ります。}

% =========================================
% 章3 典型ワークフロー
% =========================================
\chap{3}{典型ワークフロー（6ステップ）}

実際の操作の流れを、田中 太郎さんを例に並べます。すべてアプリの中で完結します。

\vspace{0.4em}
\begin{center}\sffamily\small
\begin{tikzpicture}[
  node distance=0.55cm,
  step/.style={rounded corners=4pt,draw=ipBlue!30,fill=white,minimum width=2.4cm,minimum height=1.2cm,align=center,inner sep=4pt,font=\sffamily\scriptsize,line width=0.6pt,drop fuzzy shadow=ipNavy!10},
  hi/.style={rounded corners=4pt,fill=ipOrange!10,draw=ipOrange,minimum width=2.4cm,minimum height=1.2cm,align=center,inner sep=4pt,font=\sffamily\scriptsize\bfseries,line width=0.8pt,drop fuzzy shadow=ipOrange!20},
  arr/.style={-{Stealth[length=2.2mm]},line width=0.9pt,ipBlue!60},
]
\node[step] (s1) {\textbf{1. 教材を作る}\\\scriptsize/materials\\\textcolor{ipGray}{\faBook}};
\node[step,right=of s1] (s2) {\textbf{2. 生徒を選ぶ}\\\scriptsize/students\\\textcolor{ipGray}{\faUser}};
\node[step,right=of s2] (s3) {\textbf{3. 割り当てる}\\\scriptsize tab=materials\\\textcolor{ipGray}{\faPlus}};
\node[hi,below=of s3] (s4) {\textbf{4. 結果を入れる}\\\scriptsize tab=mastery\\\textcolor{ipOrange}{\faClipboardCheck}};
\node[step,left=of s4] (s5) {\textbf{5. 反映}\\\scriptsize Ctrl+S\\\textcolor{ipGray}{\faSave}};
\node[step,left=of s5] (s6) {\textbf{6. 個別印刷}\\\scriptsize/print\\\textcolor{ipGray}{\faPrint}};

\draw[arr] (s1) -- (s2);
\draw[arr] (s2) -- (s3);
\draw[arr] (s3.south) -- ++(0,-0.3) -| (s4.north);
\draw[arr] (s4) -- (s5);
\draw[arr] (s5) -- (s6);
\end{tikzpicture}
\end{center}

\vspace{0.3em}
\begin{tipbox}
\textbf{\textcolor{ipOrange}{\faLightbulb}\ ポイント}\\[2pt]
ステップ\,\textbf{4}（結果を入れる）と\,\textbf{5}（反映）の間で\textbf{進捗が次の章へ自動で進む}のが、このアプリの一番の特徴です。
\end{tipbox}

\subsection*{ステップごとの実際の操作}

\begin{actionbox}
\textbf{\faBook\ \ 1. 教材を作る}\,\arrowto\,\texttt{/materials}\\
ヘッダ右の \btn{+ 教材追加}\,\arrowto\,ダイアログで\textbf{教材名}と\textbf{教科}（数学・英語・国語・理科・社会の色つきチップから選択）を入力\,\arrowto\,\btn{登録する}\,。\\
作成したカードの\,「\textbf{範囲を管理}」を開いて\,\btn{+ 範囲を追加}\,\arrowto\,章タイトル・問題PDF・解答PDFをドラッグ\&ドロップ。
\end{actionbox}

\begin{actionbox}
\textbf{\faUser\ \ 2. 生徒を選ぶ}\,\arrowto\,\texttt{/students}\\
画面上部の\textbf{生徒切替セレクタ}（プルダウン）から対象の生徒を選ぶ。これだけでその生徒の詳細パネルに切り替わります。
\end{actionbox}

\begin{actionbox}
\textbf{\faPlus\ \ 3. 教材を割り当てる}\,\arrowto\,\texttt{?tab=materials}\\
「\textbf{割り当て管理}」タブを開く\,\arrowto\,下半分の\textbf{緑枠「追加可能な教材」}カード一覧から、いま作った教材をクリック\,\arrowto\,上半分の\textbf{青枠「割り当て中」}に移動し、円形プログレス0\%＋ステッパー\,「1/N」\,でカードが表示される。
\end{actionbox}

\begin{actionbox}
\textbf{\faClipboardCheck\ \ 4. 結果を入れる}\,\arrowto\,\texttt{?tab=mastery}\\
「\textbf{定着度入力}」タブを開く\,\arrowto\,ツールバー左の\textbf{講師セレクト}から担当講師を選ぶ（未選択だと反映できない）\,\arrowto\,表（教材名 / 現在の範囲 / 未実施 / 得点 / 満点 / 合格 / 次回の範囲）の\textbf{得点・満点}を入れて\textbf{合格}にチェック。\\
矢印キーで隣のセルへ移動、Space で合格切替もできます。
\end{actionbox}

\begin{actionbox}
\textbf{\faSave\ \ 5. 反映（Ctrl+S）}\\
ツールバー右の \btnG{反映 (Ctrl+S)}\,\arrowto\,「採点お疲れ様でした。返却お願いします。」のトーストが出て、
\begin{itemize}[leftmargin=1.4em,itemsep=1pt]
  \item 合格した教材の\textbf{現在のポインタが +1}（次の章へ）
  \item 合格しなかった教材は\textbf{再実施扱い}（連続2回でホームに赤フラグ）
\end{itemize}
が同時に走ります。\textbf{割り当て管理タブに戻る}と、円形プログレスが新しい\%に更新されています。
\end{actionbox}

\begin{actionbox}
\textbf{\faPrint\ \ 6. 個別印刷}\,\arrowto\,\texttt{/print}\\
キュー一覧は\textbf{生徒ごとにグループ化}されています。対象の生徒のカードを開く\,\arrowto\,\textbf{問題セクション}（青）と\textbf{解答セクション}（黄）に分かれた行を確認\,\arrowto\,カード右上の \btnO{印刷}\,ボタン。\\
\arrowto\,その生徒の項目だけをまとめて\textbf{1つのPDFに結合}\,\arrowto\,新しいタブで開いて自動印刷ダイアログ。
\end{actionbox}

% =========================================
% 章4 画面ごとの詳細
% =========================================
\chap{4}{画面ごとの使い方（詳細）}

ここから「\textbf{この画面のここを押すと、こうなる}」を、画面ごとに並べていきます。

\vspace{0.5em}

% --- ダッシュボード ---
\begin{screenbox}{\faTachometerAlt\ \  ダッシュボード（ログインした最初に開く）}
\textbf{塾全体の様子をひと目で。要注意の生徒には自動で印が立ちます。}

\begin{actionbox}
\textbf{画面の上の方：}3つの数字カードが並びます。\\[2pt]
\hspace*{1em}\textcolor{ipBlue}{■}\,\textbf{生徒数}\quad
\textcolor{ipBlue}{■}\,\textbf{教材数}\quad
\textcolor{ipBlue}{■}\,\textbf{今週の学習回数}（先週比トレンドつき）
\end{actionbox}

\begin{actionbox}
\textbf{要注意の生徒が自動で出ます（2種類）：}
\begin{itemize}[leftmargin=1.4em,itemsep=2pt]
  \item \textcolor{ipOrange}{\faExclamationCircle\ \textbf{完了間近リマインド}}：あと1〜2回で教材が終わる生徒。
  \item \textcolor{ipRed}{\faExclamationTriangle\ \textbf{正答率低下リマインド}}：直近2回続けて60点未満の単元。
\end{itemize}
\textbf{右のチェックボタン}を押す\arrowto「対処済み」になって一覧から消えます（取消もワンクリック）。\\
\textbf{生徒の名前をクリック}\arrowto その子の詳細画面（教材タブ／分析タブ）にジャンプ。
\end{actionbox}

\begin{actionbox}
\textbf{下にグラフ4枚＋進捗マトリクス表：}
\begin{itemize}[leftmargin=1.4em,itemsep=1pt]
  \item 生徒ごとの平均進捗率（横棒）
  \item 過去8週間の学習量の推移（折れ線）
  \item 進捗率の分布（0-25\%／26-50\%／51-75\%／76-100\%）
  \item 生徒×教材の進捗マトリクス表（誰がどの教材をどこまで進んでいるか）
\end{itemize}
\end{actionbox}
\end{screenbox}

\vspace{0.6em}

% --- 生徒 ---
\begin{screenbox}{\faUserGraduate\ \  生徒画面}
\textbf{左に生徒リスト、右にその子の詳細。タブで内容を切り替えます。}

\begin{actionbox}
\textbf{左の \btn{+} ボタン}\arrowto 新しい生徒を追加するダイアログ（名前・学年など）。\\
\textbf{右側のタブは4枚：}\,\,
\tab{\faClipboardCheck\ 定着度入力}\,\,\tab{\faBook\ 割り当て管理}\,\,\tab{\faChartBar\ 分析}\,\,\tab{\faUsers\ 生徒一覧}
\end{actionbox}

\begin{actionbox}
\textbf{\tab{\faClipboardCheck\ 定着度入力}（このアプリの心臓部）：}\\
エクセルのような表形式で、教材ごとの「点数 / 満点」「合格チェック」を直接入力。\\
講師セレクト\,$\to$\,各セルに数値\,$\to$\,合格にチェック\,$\to$\,\btn{反映 (Ctrl+S)}\\
反映と同時に\textbf{合格した教材は次の章へ自動で進み}、不合格は再実施扱いになります。\\
2回続けて60点未満だと、ダッシュボードに\textcolor{ipRed}{\textbf{赤い警告}}が立ちます。
\end{actionbox}

\begin{actionbox}
\textbf{\tab{\faBook\ 割り当て管理}：}
教材ごとに\textbf{円形プログレス}（進捗\%）と\,\btn{−}\,/\,\btn{+}\,の章送りボタン。\\
ノードをクリックすると、その章のPDF（問題／解答／復習用）を画面内で\textbf{プレビュー}できます。\\
未割り当ての教材は下の「追加可能な教材」から\textbf{ワンクリックで割り当て}。試験教材も同様。
\end{actionbox}

\begin{actionbox}
\textbf{\tab{\faChartBar\ 分析}：}
\begin{itemize}[leftmargin=1.4em,itemsep=1pt]
  \item 正答率の推移（折れ線）
  \item 学習ペース（週あたりこなした単元の数）
  \item 進捗タイムライン（いつどの単元を完了したか）
\end{itemize}
\end{actionbox}
\end{screenbox}

\vspace{0.6em}

% --- 教材 ---
\begin{screenbox}{\faBook\ \  教材管理画面（最初の登録のときに使います）}
\textbf{教材を「章・回」単位で登録する画面です。}

\begin{actionbox}
\textbf{科目ごとに分けて表示（5教科で色分け）：}\\
\tab{数学（青）}\,\tab{英語（赤）}\,\tab{国語（緑）}\,\tab{理科（紫）}\,\tab{社会（橙）}\,\tab{その他}\\[2pt]
科目の名前をクリックで開閉。検索バーで教材名から絞り込み。
\end{actionbox}

\begin{actionbox}
\textbf{各教材カードの \btn{+}}\arrowto 章・回の追加。\\
\textbf{各章ごとに4種類のPDFを紐付けられます：}
\begin{itemize}[leftmargin=1.4em,itemsep=1pt]
  \item 問題PDF
  \item 解答PDF
  \item 復習用 問題PDF
  \item 復習用 解答PDF
\end{itemize}
\textbf{ドラッグ\&ドロップでアップロード}できます。両面印刷フラグも各章ごとに設定可。
\end{actionbox}

\begin{tipbox}
\textbf{\textcolor{ipOrange}{\faLightbulb}\ ポイント}\\
教材カード上の\textbf{小さなドット}でPDFの揃い具合が一目で分かります（緑＝あり／グレー＝なし）。
\end{tipbox}
\end{screenbox}

\vspace{0.6em}

% --- 単語テスト ---
\begin{screenbox}{\faSpellCheck\ \  単語テスト画面}
\textbf{単語帳の管理。エクセルからコピペで取り込めます。}

\begin{actionbox}
\btn{+ 新規単語帳}\,\arrowto 単語帳の名前・科目を入れて作成。\\
\btn{CSV取り込み}\,\arrowto エクセルの内容を貼付け、またはCSVファイル選択。\\
\arrowto \textbf{先頭の見出し（番号／単語／意味）を自動で読み取り}\arrowto 確認画面\arrowto 確定でまとめて登録。
\end{actionbox}

\begin{actionbox}
クリックで単語帳の中身（番号・単語・意味の3列）を確認。\\
教材と紐付けておけば、後で「\textbf{この範囲だけの単語テスト}」を作る土台にもなります。
\end{actionbox}
\end{screenbox}

\vspace{0.6em}

% --- 試験管理 ---
\begin{screenbox}{\faChartBar\ \  試験管理画面（合格力の見える化）}
\textbf{共通テスト・大学別過去問の管理と、成績の多角的な分析。}

\begin{actionbox}
\textbf{2タブ構成：} \tab{分析}\,\,\tab{試験教材}
\end{actionbox}

\begin{actionbox}
\textbf{\tab{試験教材}：}
共テ・模試／大学過去問のサブタブ。各試験カードに年度・大学・学部・教科（満点）を登録。
カードを開いて科目の追加／削除、生徒への割り当てができます。
\end{actionbox}

\begin{actionbox}
\textbf{\tab{分析}（このアプリの目玉）：}\\
生徒セレクト×試験セレクトの組み合わせで\textbf{8種類のグラフ}が出ます：
\begin{itemize}[leftmargin=1.4em,itemsep=1pt]
  \item 教科別スコア（横棒）
  \item 教科バランス（\textbf{レーダー}：実績／目標／クラス平均を重ね表示）
  \item 個人 vs クラス平均（教科別比較）
  \item 目標との差
  \item 得点率の推移（複数試験ぶん折れ線）
  \item 全体の点数分布（ヒストグラム）
  \item クラスの教科別平均
  \item \textbf{圧縮スコア計算機}（次のヒント参照）
\end{itemize}
\end{actionbox}

\begin{tipbox}
\textbf{\textcolor{ipOrange}{\faStar}\ 目玉：圧縮スコア計算}\\
大学・学部ごとに\textbf{配点の重み}を登録しておけば、生徒の得点を志望大学の配点に\textbf{自動で換算}。\\
「いまの成績で、志望大学のボーダーまでどのくらいか」が、即座に出ます。
\end{tipbox}
\end{screenbox}

\vspace{0.6em}

% --- 印刷 ---
\begin{screenbox}{\faPrint\ \  印刷画面（生徒ごとに、まとめて出力）}
\textbf{印刷は生徒1人ずつ、その子の必要分だけ出せます。}

\begin{actionbox}
\textbf{2タブ構成：}\tab{\faListUl\ キュー}\,\,\tab{\faHistory\ ジョブ履歴}
\end{actionbox}

\begin{actionbox}
\textbf{\tab{キュー}：}印刷待ちの一覧が\textbf{生徒ごとにグループ化}されて並びます。\\
各生徒のカードは「\textcolor{ipBlue}{問題}」セクションと「\textcolor{ipOrange}{解答}」セクションに分かれていて、
それぞれにプレビュー／削除のボタンがついています。\\
カード右上の \btnG{印刷}\,ボタンを押すと、その生徒のぶんを\textbf{1つのPDFに結合}\arrowto 新しいタブで開いて自動印刷。
\end{actionbox}

\begin{actionbox}
\textbf{ツールバーの便利機能：}
\begin{itemize}[leftmargin=1.4em,itemsep=1pt]
  \item \btn{手動追加}\,：生徒×教材×章を選んでキューに追加
  \item \btn{全生徒の次回分を自動追加}\,：一気に全員ぶん並べる
  \item \btn{元に戻す}\,：直前の印刷ジョブをキューに戻して再実行
  \item 表示切替：「問題＋解答／問題のみ／解答のみ」
\end{itemize}
\end{actionbox}

\begin{actionbox}
\textbf{\tab{ジョブ履歴}：}過去の印刷ジョブの状態（受付中／送信中／完了／失敗）と、不足PDF件数を確認できます。
\end{actionbox}
\end{screenbox}

\vspace{0.6em}

% --- PDF / 講師 / 設定 ---
\begin{screenbox}{\faFilePdf\ PDF置き場 ／ \faChalkboardTeacher\ 講師 ／ \faUserShield\ アカウント管理}
\textbf{細かい設定画面たち。普段は触らず、初期設定のときに使う場所です。}

\begin{actionbox}
\textbf{\faFilePdf\ \ PDF置き場：}
教材PDFをフォルダごとに一覧／プレビュー／アップロード／削除。ドラッグ\&ドロップで一気にアップロード可。
\end{actionbox}

\begin{actionbox}
\textbf{\faChalkboardTeacher\ \ 講師：}
講師の名前を登録します。\textbf{誰が教えた回か}を定着度の記録と紐付ける用。
\end{actionbox}

\begin{actionbox}
\textbf{\faUserShield\ \ アカウント管理：}
画面に入れる人を管理。役割は2種類。
\begin{itemize}[leftmargin=1.4em,itemsep=1pt]
  \item \tab{\faCrown\ 塾長（管理者）}：すべての画面OK
  \item \tab{\faUser\ 先生（トレーナー）}：生徒・進捗の確認のみOK（教材編集・印刷・試験管理は不可）
\end{itemize}
追加するときは\textbf{Googleアカウントのメールアドレス}と表示名・役割を入れて \btn{追加}\,。
\end{actionbox}
\end{screenbox}

% =========================================
% 章4 自動化のキモ
% =========================================
\chap{5}{自動化のキモ\,3 つ}

ここがこのアプリ最大の「楽になるポイント」です。

\vspace{0.4em}
\begin{tcolorbox}[enhanced,breakable,colback=ipMint,colframe=ipGreen,boxrule=0.8pt,arc=5pt,top=14pt,bottom=12pt,left=14pt,right=14pt,
title={\sffamily\bfseries\large\faMagic\ \  1.\,\,点数を入れたら、進捗が勝手に進む},coltitle=white,colbacktitle=ipGreen,
attach boxed title to top left={xshift=14pt,yshift=-12pt},
boxed title style={enhanced,arc=4pt,boxrule=0pt,top=3pt,bottom=3pt,left=10pt,right=10pt,colback=ipGreen,interior style={fill=ipGreen}}]
\sffamily
\vspace{0.4em}
定着度入力で「点数」「満点」を入れて\textbf{合格}にチェック\arrowto\,\btn{反映 (Ctrl+S)}\,を押すと、
\begin{itemize}[leftmargin=1.4em,itemsep=2pt]
  \item 合格した教材は\textbf{次の章へ自動で進み}
  \item 不合格は\textbf{再実施扱い}（連続2回で赤フラグ）
  \item 該当ノードが\textbf{印刷キューに自動追加}される（運用設定により）
\end{itemize}
1人の生徒で\textbf{5教科ぶんを一度にまとめて}入れられ、Ctrl+S 1回で全部反映。
\end{tcolorbox}

\vspace{0.5em}
\begin{tcolorbox}[enhanced,breakable,colback=ipBlue!6,colframe=ipBlue,boxrule=0.8pt,arc=5pt,top=14pt,bottom=12pt,left=14pt,right=14pt,
title={\sffamily\bfseries\large\faExclamationTriangle\ \  2.\,\,要注意の生徒に、自動で印が立つ},coltitle=white,colbacktitle=ipBlue,
attach boxed title to top left={xshift=14pt,yshift=-12pt},
boxed title style={enhanced,arc=4pt,boxrule=0pt,top=3pt,bottom=3pt,left=10pt,right=10pt,colback=ipBlue,interior style={fill=ipBlue}}]
\sffamily
\vspace{0.4em}
ダッシュボードを開くと、\textbf{先生がチェックしなくても}次の生徒に印が立っています。
\begin{itemize}[leftmargin=1.4em,itemsep=2pt]
  \item \textcolor{ipOrange}{\textbf{完了間近}}：あと1〜2回で教材が終わる生徒（次の教材を準備するヒント）
  \item \textcolor{ipRed}{\textbf{正答率低下}}：直近2回続けて60点未満（指導の見直しサイン）
\end{itemize}
\textbf{対処済み}にしたら一覧から消え、間違ってチェックした場合は\textbf{ワンクリックで取り消し}できます。
\end{tcolorbox}

\vspace{0.5em}
\begin{tcolorbox}[enhanced,breakable,colback=ipPurple!8,colframe=ipPurple,boxrule=0.8pt,arc=5pt,top=14pt,bottom=12pt,left=14pt,right=14pt,
title={\sffamily\bfseries\large\faCalculator\ \  3.\,\,志望大学の配点で、合格力を換算},coltitle=white,colbacktitle=ipPurple,
attach boxed title to top left={xshift=14pt,yshift=-12pt},
boxed title style={enhanced,arc=4pt,boxrule=0pt,top=3pt,bottom=3pt,left=10pt,right=10pt,colback=ipPurple,interior style={fill=ipPurple}}]
\sffamily
\vspace{0.4em}
試験管理の分析タブで\textbf{圧縮スコア計算機}を起動。\\
あらかじめ「東大 理科一類」「京大 工学部」のような\textbf{大学・学部ごとの配点表}を登録しておけば、
\begin{itemize}[leftmargin=1.4em,itemsep=2pt]
  \item いまの模試成績を\textbf{志望大学の配点に自動換算}
  \item ボーダーまでの距離が\textbf{即座に表示}
  \item 「東大」「京大」など複数候補を\textbf{ワンクリックで切り替え}
\end{itemize}
\end{tcolorbox}

% =========================================
% 章5 裏で覚えていること
% =========================================
\chap{6}{裏で何を覚えているか}

塾長が中身を細かく追う必要はありません。「これだけは記録に残ってます」というポイントだけ。

\vspace{0.4em}
\begin{center}\sffamily\small
\renewcommand{\arraystretch}{1.4}
\begin{tabularx}{\linewidth}{@{}>{\bfseries\color{ipNavy}}l X@{}}
\toprule
覚えている情報 & 何のために \\
\midrule
誰がどの教材のどこにいるか & 「次にやる範囲」を即決められる \\
1回ごとの点数・正答率・日付・担当の先生 & 弱点と先生ごとの傾向が見える \\
進捗の履歴（前進・戻し・印刷など） & ダッシュボードのトレンドグラフの元データ \\
模試の点数・目標・大学別配点表 & 目標達成度／圧縮スコアの自動計算 \\
印刷キュー・実行履歴・プリンタ設定 & 生徒ごとの印刷とトラブル時の取り消し \\
教材ごとの問題／解答／復習用PDF（4種類） & その場プレビュー＋まとめて印刷 \\
完了したものの倉庫／対処済み印 & 「終わったもの」をきれいに隠す \\
\bottomrule
\end{tabularx}
\end{center}

\vspace{0.3em}
\textcolor{ipGray}{\small※ ログイン権限は \textbf{塾長} と \textbf{先生} の2段階。間違えて教材を弄られないよう、先生アカウントは触れる範囲を絞っています。}

% =========================================
% 章6 進捗
% =========================================
\needspace{20\baselineskip}
\chap{7}{今どこまでできているか}

体感では \textbf{普段使う機能はだいたい完成}、いまは細かい見た目を整えているところです。

\vspace{0.4em}
\begin{center}\sffamily\small
\renewcommand{\arraystretch}{1.4}
\begin{tabularx}{\linewidth}{@{}>{\bfseries\color{ipNavy}}l c c X@{}}
\toprule
こんな機能 & 完成度 & ゲージ & 状況 \\
\midrule
生徒・教材・進捗の管理 & \textcolor{ipGreen}{\textbf{ほぼ完成}} & \progbar{95}{ipGreen} & 日常運用は十分まわせます \\
定着度入力（表形式UI） & \textcolor{ipGreen}{\textbf{ほぼ完成}} & \progbar{92}{ipGreen} & Ctrl+S 一括反映・自動進行まで動作 \\
要注意の自動検出（リマインド） & \textcolor{ipGreen}{\textbf{ほぼ完成}} & \progbar{90}{ipGreen} & 対処済みフラグ・取消も対応 \\
印刷まわり全般 & \textcolor{ipGreen}{\textbf{ほぼ完成}} & \progbar{90}{ipGreen} & 生徒個別印刷・取り消し対応 \\
試験分析（8種グラフ・圧縮スコア） & \textcolor{ipGreen}{\textbf{ほぼ完成}} & \progbar{88}{ipGreen} & 大学別配点換算まで対応 \\
単語テスト & \textcolor{ipBlue}{\textbf{動いている}} & \progbar{75}{ipBlue} & 取り込みOK／出題機能は今後 \\
画面の細かい整え & \textcolor{ipOrange}{\textbf{調整中}} & \progbar{70}{ipOrange} & 行間・崩れ・文字化けの直し \\
本番運用へ向けた整備 & \textcolor{ipOrange}{\textbf{これから}} & \progbar{40}{ipOrange} & 自動チェック・運用説明書 \\
\bottomrule
\end{tabularx}
\end{center}

\vspace{0.5em}
\subsection*{\faClock\ \ 最近やっていたこと（直近の作業から抜粋）}
\begin{itemize}[leftmargin=1.4em,itemsep=2pt]
  \item 答案用紙まわりの表示の整え／印刷時の崩れ修正
  \item 表の行の高さ・列幅のちょうどいい値さがし
  \item 単語テストの一括取り込み（見出しの自動判定・文字化けの対策）
  \item 一気に保存しても止まらないよう改善
  \item ログイン権限まわりの整備
  \item 見える化・分析の強化／生徒追加画面／ログまわりの整理
\end{itemize}

% =========================================
% 章7 次にやれること
% =========================================
\needspace{18\baselineskip}
\chap{8}{次にやれること（塾長と相談したいリスト）}

「重さの目安」つきで並べました。「中」くらいまでは手の届く範囲です。

\vspace{0.4em}
\begin{center}\sffamily\small
\renewcommand{\arraystretch}{1.4}
\begin{tabularx}{\linewidth}{@{}>{\bfseries\color{ipNavy}}l c X@{}}
\toprule
やりたいこと & 重さ & 効果 \\
\midrule
単語テストの「出題→提出→自動採点」画面 & \tab{中} & 単語確認の時間がグッと減る \\
保護者・生徒向けの進捗共有ページ（読むだけ） & \tab{中} & 通塾モチベの可視化 \\
模試結果の取り込み（ファイル／写真OCR） & \tab{中〜大} & 成績入力の手間がなくなる \\
LINE／Slack への通知連携 & \tab{小} & 要注意フラグを自動でお知らせ \\
夜のうちに翌日ぶん印刷を自動準備 & \tab{中} & 朝イチでもう紙が積まれている \\
本番運用に向けた自動チェック・バックアップ & \tab{中} & 安心感（裏側の整備） \\
\bottomrule
\end{tabularx}
\end{center}

% =========================================
% 章8 おわりに
% =========================================
\chap{9}{おわりに}

\begin{tcolorbox}[enhanced,colback=ipNavy!4,colframe=ipNavy,boxrule=0.8pt,arc=5pt,top=14pt,bottom=14pt,left=18pt,right=18pt]
\sffamily
{\fontsize{13}{18}\selectfont\textbf{\textcolor{ipNavy}{\faQuoteLeft}\ \ まとめ}}\\[8pt]

\begin{tcolorbox}[colback=ipSoft,colframe=ipOrange!50,boxrule=0.6pt,arc=3pt,
left=12pt,right=12pt,top=8pt,bottom=8pt,
borderline west={2.4pt}{0pt}{ipOrange}]
\sffamily\bfseries\textcolor{ipNavy}{ダッシュボードを開けば塾の今が見え、点数を入れたら進捗が進み、
模試を入れたら合格力が見え、生徒ごとに必要な紙だけ印刷できる。}
\end{tcolorbox}

\vspace{8pt}
\textcolor{ipNavy}{これが今の iPlus Sys の立ち位置です。生徒・教材・定着度・模試・印刷まで、
\textbf{塾の運用に必要なものを全部このアプリの中で}巻き取れる状態になっています。}\par
\vspace{6pt}
\textcolor{ipGray}{細かい挙動や追加してほしい画面があれば、いつでもこの資料に追記する形でお伝えください。}
\end{tcolorbox}

\end{document}
